\chapter{The Krylov criterion and generic Hohenberg--Kohn}
\label{chap:generic}

\section{Krylov polynomials}

\begin{definition}[Krylov matrix]
  \label{def:kry}
  \lean{DFT.kry}
  \leanok
  For a matrix $H$ on $\Conf$, a configuration $c$ and an enumeration $\mathrm{ord}:\Conf\to\N$, the matrix
  $\Kry(H,c)$ has columns $H^{\mathrm{ord}(j)}e_c$.
\end{definition}

\begin{lemma}[Vanishing amplitudes make the Krylov matrix singular]
  \label{lem:det_kry_eq_zero}
  \lean{DFT.det_kry_eq_zero}
  \leanok
  \uses{def:kry}
  If $H$ is symmetric, $H\psi=E\psi$, $\psi\neq0$ and $\psi_c=0$, then $\det\Kry(H,c)=0$.
\end{lemma}
\begin{proof}
  \leanok
  $\ip{\psi}{H^ke_c}=E^k\psi_c=0$, so $\psi$ is a nonzero vector in the left kernel.
\end{proof}

\begin{definition}[Krylov polynomial]
  \label{def:krylovPoly}
  \lean{DFT.hamP, DFT.krylovPoly}
  \leanok
  \uses{def:kry, def:hamN}
  Treating the potential as polynomial variables,
  $P_c(v)=\det\Kry(H(v),c)\in\R[v_i : i\in V]$.
\end{definition}

\begin{lemma}
  \label{lem:eval_krylovPoly}
  \lean{DFT.hamP_eval, DFT.eval_krylovPoly}
  \leanok
  \uses{def:krylovPoly}
  Evaluating $P_c$ at $v$ gives $\det\Kry(H(v),c)$.
\end{lemma}
\begin{proof}
  \leanok
  Evaluation is a ring homomorphism; it commutes with matrix powers and determinants.
\end{proof}

\begin{definition}[Rigidity and non-unique $v$-representability]
  \label{def:NonUV}
  \lean{DFT.Rigid, DFT.NonUV}
  \leanok
  \uses{def:hamN, def:IsGroundState}
  The occupations are \emph{rigid} if $\sum_iw_in_i(c)$ constant in $c$ forces $w$ constant.
  $v$ is \emph{non-uniquely $v$-representable} if some ground state of $H(v)$ is a ground state of $H(v')$
  with $v-v'$ non-constant.
\end{definition}

\begin{theorem}
  \label{thm:nonUV_imp_krylov_zero}
  \lean{DFT.nonUV_imp_krylov_zero}
  \leanok
  \uses{def:NonUV, lem:det_kry_eq_zero, lem:eval_krylovPoly, thm:isGroundState_eigen}
  With rigid occupations, a non-uniquely $v$-representable $v$ is a zero of some Krylov polynomial $P_c$.
\end{theorem}
\begin{proof}
  \leanok
  If the shared ground state had full support, rigidity would force $v-v'$ constant. So some amplitude
  vanishes; apply Lemma~\ref{lem:det_kry_eq_zero}.
\end{proof}

\section{Zero sets of polynomials}

\begin{theorem}[Identity theorem for polynomials]
  \label{thm:zeroSet_closed_nowhere}
  \lean{DFT.eq_zero_of_vanish_on_open, DFT.zeroSet_closed_nowhere}
  \leanok
  A real polynomial vanishing on a nonempty open set is zero. Hence the zero set of a nonzero polynomial is
  closed with empty interior.
\end{theorem}
\begin{proof}
  \leanok
  Polynomials are analytic; apply the identity theorem on the connected space $\R^V$.
\end{proof}

\begin{theorem}[Generic HK, topological form]
  \label{thm:generic_hk}
  \lean{DFT.generic_hk}
  \leanok
  \uses{thm:nonUV_imp_krylov_zero, thm:zeroSet_closed_nowhere}
  Assume (H$^*$): every Krylov polynomial $P_c$ is nonzero. Then the non-uniquely $v$-representable
  potentials lie in a closed set with empty interior.
\end{theorem}
\begin{proof}
  \leanok
  They lie in the zero set of $\prod_cP_c\neq0$.
\end{proof}

\begin{theorem}[Zero sets are null]
  \label{thm:volume_zeroSet}
  \lean{DFT.volume_zeroSet_fin, DFT.volume_zeroSet}
  \leanok
  The zero set of a nonzero real polynomial in finitely many variables has Lebesgue measure zero.
\end{theorem}
\begin{proof}
  \leanok
  Induction on the number of variables. View $P$ as a polynomial in $x_0$ whose leading coefficient $q$ is
  a nonzero polynomial in the remaining variables. Off the null set $\{q=0\}$ each slice has finitely many
  roots, so Fubini gives measure zero.
\end{proof}

\begin{theorem}[Generic HK, measure form]
  \label{thm:generic_hk_ae}
  \lean{DFT.generic_hk_ae}
  \leanok
  \uses{thm:generic_hk, thm:volume_zeroSet}
  Under (H$^*$), the non-uniquely $v$-representable potentials lie in a closed, nowhere dense, Lebesgue-null set.
\end{theorem}
\begin{proof}
  \leanok
  Combine Theorems~\ref{thm:generic_hk} and~\ref{thm:volume_zeroSet}.
\end{proof}

\section{Simultaneous non-vanishing}

\begin{definition}[Rational on generic points]
  \label{def:RatOn}
  \lean{DFT.RatOn}
  \leanok
  $f:\R^M\to\R$ is \emph{rational on generic points} if there are polynomials $P,Q$ with $Q(\theta)\ne0$ and
  $f(\theta)Q(\theta)=P(\theta)$ for every injective $\theta$.
\end{definition}

\begin{lemma}[Closure properties]
  \label{lem:RatOn_closure}
  \lean{DFT.RatOn.const, DFT.RatOn.add, DFT.RatOn.mul, DFT.RatOn.div_sub, DFT.RatOn.sum, DFT.RatOn.prod, DFT.RatOn.ite, DFT.RatOn.congr}
  \leanok
  \uses{def:RatOn}
  The class is closed under constants, sums, products, finite sums and products, case distinctions with
  constant conditions, pointwise equality on injective points, and division by $\theta_i-\theta_j$.
\end{lemma}
\begin{proof}
  \leanok
  Clear denominators; for $i=j$ the quotient is $0$ by convention.
\end{proof}

\begin{lemma}[Discriminant]
  \label{lem:discPoly}
  \lean{DFT.discPoly, DFT.discPoly_ne_zero, DFT.injective_of_disc}
  \leanok
  $D=\prod_{i\neq j}(X_i-X_j)$ is a nonzero polynomial, and $D(\theta)\neq0$ implies that $\theta$ is injective.
\end{lemma}
\begin{proof}
  \leanok
  Each factor is nonzero, as evaluation at a unit vector shows.
\end{proof}

\begin{theorem}[Simultaneous non-vanishing]
  \label{thm:exists_common_ne}
  \lean{DFT.exists_common_ne}
  \leanok
  \uses{lem:RatOn_closure, lem:discPoly}
  Finitely many functions that are rational on generic points, each nonzero at some injective point, are
  simultaneously nonzero at a common injective point. That point can moreover avoid the zero sets of
  finitely many nonzero polynomials.
\end{theorem}
\begin{proof}
  \leanok
  The numerators $P_k$ are nonzero polynomials. The product $D\prod_kP_k\prod R$ is nonzero, so it is
  nonzero at some point.
\end{proof}

\section{A certificate: two fermions on the square}

\begin{theorem}[(H$^*$) for the square]
  \label{thm:square_generic_hk}
  \lean{DFT.vsq, DFT.Hsq1, DFT.hamN_vsq, DFT.Hsq2, DFT.Hsq_pow2, DFT.Hsq3, DFT.Hsq_pow3, DFT.Hsq4, DFT.Hsq_pow4, DFT.Hsq5, DFT.Hsq_pow5, DFT.pow_cases, DFT.Ksq0, DFT.kry_sq0, DFT.det_Ksq0, DFT.Ksq1, DFT.kry_sq1, DFT.det_Ksq1, DFT.Ksq2, DFT.kry_sq2, DFT.det_Ksq2, DFT.Ksq3, DFT.kry_sq3, DFT.det_Ksq3, DFT.Ksq4, DFT.kry_sq4, DFT.det_Ksq4, DFT.Ksq5, DFT.kry_sq5, DFT.det_Ksq5, DFT.rigid_occ4, DFT.hstar_sq, DFT.square_generic_hk, DFT.square_generic_hk_ae}
  \leanok
  \uses{thm:generic_hk_ae, thm:hk_fails_square, lem:Qsq}
  For two fermions on the square, every Krylov polynomial is nonzero. Hence HK holds for Lebesgue-almost
  every potential.
\end{theorem}
\begin{proof}
  \leanok
  Lean evaluates the six $6\times6$ Krylov determinants exactly at $v=(-3,-3,-2,0)$, obtaining
  $64,-432,-448,-448,-432,64$.
\end{proof}
